\documentclass[11pt]{article}
\usepackage[margin=1in]{geometry}
\usepackage{amsmath,amssymb,booktabs,enumitem,xcolor,hyperref,listings}
\hypersetup{colorlinks=true,linkcolor=blue,urlcolor=blue}
\lstset{basicstyle=\ttfamily\small,columns=fullflexible,frame=single}
\title{DiffSim Learning Guide\\A1--A3: Verification, Boundary Conditions, and Shifted Geometry}
\author{Study notes from the DiffSim tutorials}
\date{July 2026}

\begin{document}
\maketitle
\tableofcontents
\newpage

\section{How to read the DiffSim code}

The tutorials combine mathematics, finite elements, sparse linear algebra,
and GPU kernels. The basic computational flow is always
\[
 \text{exact solution} \to \text{mesh} \to \text{assemble } A,b
 \to \text{apply boundary conditions} \to \text{solve} \to \text{measure error}.
\]

The code uses a hybrid CPU/GPU design. Normal Python, NumPy, SciPy, sparse
matrix conversion, and \texttt{splu} run on the CPU (the \emph{host}). Warp
kernels launched with \texttt{device=DEVICE} run on the GPU (the \emph{device}).
The GPU does not automatically accelerate a SciPy call.

\section{Finite-element building blocks}

\subsection{From cells to a finite-element mesh}

\texttt{build\_uniform(level, dim=2)} creates square cells. These cells are
not yet a finite-element algebraic system. \texttt{build\_mesh(tree,p)}
creates global node coordinates, element connectivity, boundary-node flags,
and polynomial-order information. For $p=1$, a square has corner nodes. For
$p=2$, it also has edge and interior interpolation points.

If an element has local nodes $a$ and basis functions $N_a$, then
\[
 u_h(x)=\sum_a N_a(x)u_a.
\]
Assembly combines element matrices into a global sparse matrix.

\subsection{Why constraints are needed}

Adaptive meshes can have a fine element next to a coarse element. A node on
the fine side may not exist on the coarse side; this is a \emph{hanging node}.
Continuity requires its value to be interpolated from independent nodes. The
constraint object stores
\[
 u_{\rm all}=T u_{\rm free}.
\]
The solver works with the independent values and reconstructs all values by
\texttt{u\_all = T @ u\_free}. On a uniform mesh there are no hanging nodes,
so $T$ is the identity.

For quadratic elements the interpolation is not generally a $0.5/0.5$ rule.
The code evaluates the appropriate Lagrange basis functions. For nodes
$\xi=-1,0,1$, evaluation at $\xi=-0.5$ gives weights
\[
 (0.375,\;0.75,\;-0.125).
\]

\section{A1: manufactured solutions and convergence}

\subsection{The PDE and MMS idea}

The model problem is Poisson's equation:
\[
 -\Delta u=f \quad \text{in }\Omega,
 \qquad u=g \quad \text{on }\partial\Omega.
\]
The method of manufactured solutions chooses a known answer $u^*$, then
defines the data so that $u^*$ solves the PDE:
\[
 f=-\Delta u^*, \qquad g=u^*|_{\partial\Omega}.
\]
For a smooth solution and linear elements, the expected $L^2$ error is
\[
 \|u-u_h\|_{L^2}=O(h^2).
\]
For quadratic elements it is typically $O(h^3)$.

The measured order between two levels is
\[
 r=\log_2\left(\frac{E_h}{E_{h/2}}\right).
\]
An order of 2 means that halving $h$ reduces the error by about four.

\subsection{A1 results}

The $p=1$ run produced errors approximately
\[
 1.302\times10^{-3},\quad 3.255\times10^{-4},\quad
 8.138\times10^{-5},
\]
with orders $2.00,2.00$. This verifies the expected second-order behavior.

The quadratic patch test used $u^*=x^2+y^2$. A quadratic finite-element
space represents this function exactly, so the error was near machine
precision ($10^{-15}$). This is not ordinary asymptotic convergence; it is a
discrete-space exactness test.

\section{Boundary conditions}

\subsection{Dirichlet row replacement}

The discrete system has the form
\[
 A u=b.
\]
If a boundary degree of freedom must equal $g_i$, one simple method replaces
its row by an identity row:
\[
 A_{i,:}\leftarrow e_i^T, \qquad b_i\leftarrow g_i.
\]
The resulting equation is exactly $u_i=g_i$.

This is easy to understand and useful in tutorials. Production codes often
eliminate known boundary degrees of freedom. Splitting unknowns into interior
and boundary parts gives
\[
 A_{II}u_I=b_I-A_{IB}g.
\]
This reduced system is smaller and can preserve symmetry when rows and
columns are treated consistently.

\subsection{Natural Neumann conditions and pinning}

After integration by parts,
\[
 \int_\Omega \nabla w\cdot\nabla u\,d\Omega
 -\int_{\partial\Omega}w\,\nabla u\cdot n\,d\Gamma
 =\int_\Omega wf\,d\Omega.
\]
If no boundary term is added on a face, the weak form naturally imposes
\[
 \nabla u\cdot n=0.
\]
With pure Neumann conditions, adding a constant to a solution gives another
solution. The matrix is singular. Pinning one node, for example $u(0,0)=1$,
removes this arbitrary constant (the gauge). It does not represent a new
physical wall condition.

For the cosine manufactured solution, the normal flux is zero on every face,
so the pure-Neumann problem is compatible. The sine solution has nonzero flux
on the $y$-faces, so using zero natural flux there produces an inconsistent
problem and a large nonconvergent error.

\section{A2 results and lesson}

With $u^*=\sin(\pi x)\sin(\pi y)$, all-face Dirichlet conditions gave
second-order convergence. The x-only case gave errors near $0.47$ and order
near zero. The reason is that the omitted $y$-faces silently impose zero
flux, while
\[
 \frac{\partial u^*}{\partial y}
 =\pi\sin(\pi x)\cos(\pi y)
\]
is generally nonzero there. Refinement cannot fix a wrong boundary condition.

The A2 performance experiment timed the CPU Python row-replacement loop.
Boundary nodes doubled from 128 to 1024, while loop time grew from
$2.23\times10^{-4}$ to $1.792\times10^{-3}$ seconds. The measured exponent
was about $0.5$ with respect to total 2-D degrees of freedom, consistent with
boundary size growing like $\sqrt n$.

\section{A3: shifted boundary method}

\subsection{Exact circle and Cartesian background mesh}

A3 solves Poisson's equation inside a disk of radius $0.3$ centered at
$(0.5,0.5)$, but the mesh is a Cartesian square mesh. The mesh is not fitted
to the circle. Elements are classified as inside or outside, and exposed
faces of retained elements create a staircase \emph{surrogate boundary}.

\subsection{The lambda criterion}

For an intercepted element, let $\theta_{\rm out}$ be its estimated outside
volume fraction. The production rule keeps it when
\[
 \theta_{\rm out}\le \lambda.
\]
Thus $\lambda=0$ keeps only fully interior elements, while $\lambda=1$
keeps every element touched by the circle. Fully interior elements are always
kept and fully exterior elements are dropped.

\subsection{Projection data}

For every Gauss quadrature point $x_q$ on a surrogate face, the geometry
oracle finds a closest true-boundary point $y$. The displacement is
\[
 d=y-x_q.
\]
The true normal is the normalized level-set gradient at $y$:
\[
 n=\frac{\nabla\psi(y)}{\|\nabla\psi(y)\|},
\]
oriented outward from the physical domain. It is conceptually separate from
$d$, although an exact closest-point projection on a smooth surface makes the
two directions parallel.

The displacement tells the method \emph{where} the true boundary is. The
normal tells it \emph{which direction is perpendicular to the boundary} and
is needed for the flux $\nabla u\cdot n$ and Nitsche terms.

\subsection{Taylor shift and the Hessian}

The surrogate point is not the true boundary point. The method approximates
the value at $y=x+d$ by Taylor expansion:
\[
 Su(x)=u(x)+\nabla u(x)\cdot d
 +\frac12d^T H(u)(x)d+\cdots.
\]
The Hessian is the matrix of second derivatives:
\[
 H(u)=
 \begin{bmatrix}u_{xx}&u_{xy}\\u_{yx}&u_{yy}\end{bmatrix}.
\]
For $p=1$, the first-order shift is used. For $p=2$, the Hessian term is
needed to obtain the matching third-order $L^2$ behavior. The sign of $d$
matters: when the surrogate changes from inside to outside, $d$ reverses,
and the signed Taylor term reverses too.

\subsection{Nitsche terms}

The shifted SBM Dirichlet form contains, schematically,
\[
 -\int_{\widetilde\Gamma} w(\nabla u\cdot\widetilde n)\,dS
 -\int_{\widetilde\Gamma}(\nabla w\cdot\widetilde n)Su\,dS
 +\int_{\widetilde\Gamma}\frac{\alpha}{h}(Sw)(Su)\,dS.
\]
The first term is consistency, the second is adjoint consistency, and the
third is penalty/stability. The test function is also shifted:
\[
 Sw=w+\nabla w\cdot d+\frac12d^TH(w)d.
\]
This does not physically move the function; it approximates its value at the
true boundary while integration remains on the staircase.

\subsection{A3 measured results}

For $p=1$, the exact-circle experiment gave errors
\[
3.210\times10^{-3},\;6.615\times10^{-4},\;1.508\times10^{-4}
\]
and orders $2.28,2.13$, consistent with $O(h^2)$. For $p=2$, the errors were
\[
1.726\times10^{-4},\;1.423\times10^{-5},\;1.808\times10^{-6}
\]
and orders $3.60,2.98$, confirming the expected $O(h^3)$ rate on the finer
interval.

With $\lambda=1$, the p1 errors were
\[
9.324\times10^{-3},\;3.753\times10^{-3},\;6.083\times10^{-4},
\]
with orders $1.31,2.63$. The coarse interval is pre-asymptotic; the finer
interval is approaching the expected second-order behavior. The method works
whether the surrogate is inside or outside because $d$ is signed.

\subsection{Sampled GridSDF geometry}

\texttt{GridSDF} stores signed-distance values at a $129\times129$ grid of
corner points and uses multilinear interpolation. Negative values mean
inside, positive values mean outside, and the interpolated zero set
approximates the circle. This models geometry obtained from voxel data or a
sampled level set rather than an analytic formula.

The strict closest-point test failed at 8 of 80 points because the sampled
field is only piecewise smooth. Allowing a 10\% fallback produced errors
\[
9.317\times10^{-3},\;3.791\times10^{-3},\;6.401\times10^{-4}
\]
and orders $1.30,2.57$, close to the exact-circle $\lambda=1$ result. The
physical lesson is that geometry representation can eventually create an
error floor: once the FEM mesh is finer than the geometry data, refining the
FEM mesh alone cannot recover the exact shape.

\section{Performance and the research lesson}

For a 2-D uniform refinement, the number of degrees of freedom grows by about
four when $h$ is halved. If $t\propto n^\alpha$, then
\[
 \alpha=\frac{\log(t_2/t_1)}{\log(n_2/n_1)}.
\]
The A1 performance study found approximately linear assembly and a
superlinear sparse solve. At level 8, representative GPU-visible timings
were assembly $0.1501$ seconds and solve $0.2433$ seconds. The Warp assembly
kernels ran on the RTX GPU, but SciPy's \texttt{splu} solve remained on the
CPU. Thus the solve was the bottleneck.

The broader lesson is to measure both accuracy and cost. A method can have
the correct convergence order but still become unusable if its solve or
geometry stage scales poorly. Conversely, a fast calculation is not useful if
its boundary conditions or geometry are wrong.

\section{Final checklist}

\begin{enumerate}[leftmargin=*]
\item Derive $f$ and $g$ from a known manufactured solution.
\item Check errors and convergence orders, not just one solution plot.
\item Distinguish physical boundary conditions from mesh-continuity constraints.
\item Remember that omitted weak boundary terms impose natural zero flux.
\item Use signed displacement $d$ for shifted geometry and the true normal for flux.
\item Match Taylor-shift order to polynomial order: p1 uses first order, p2 uses Hessian correction.
\item Expect pre-asymptotic fluctuations on coarse meshes.
\item Measure scaling and identify the eventual bottleneck.
\item Treat sampled geometry as an additional numerical approximation with its own error floor.
\end{enumerate}

\end{document}
